% Conclusion

This study characterised the underestimation of vessel diameter in OCTA and evaluated how it can be corrected. We established ground truth in the in vivo mouse retina using Intralipid contrast enhancement. Binary mask based diameter calculation was then compared against model-based profile fitting, on both maximum and mean intensity projections. Every method underestimated vessel diameter before contrast. The generalised Gaussian $1/e^2$ metric retained the strongest linear relationship with the ground truth under these conditions, so the residual underestimation is systematic and can be removed by calibration. Post-contrast, the same metric applied to the mean intensity projection provided results closest to the ground truth of all combinations tested. Binary mask based measurements depended on the threshold and the contrast condition, with the post-contrast slope falling steeply as the threshold increased. The practical consequence of that dependence was demonstrated by the flicker stimulation results. As the angiogram signal declined over the course of the acquisition, binary mask based diameters followed it and reported vasoconstriction. The model-based fits were effectively independent of intensity and recovered the expected vasodilation. We therefore recommend a generalised Gaussian fit to the mean intensity projection for quantitative OCTA vessel diameter measurement, with a calibrated correction applied where contrast enhancement is not available.
